\documentclass[hyperref={pdfpagelabels=false},table]{beamer}

\input{lec_common}
\begin{document}

\part{Unit 5: Linear Algebra}

%\frame{\tableofcontents}


\begin{frame}[fragile]{The \texttt{array} type}
Lists are almost like vectors but the operations on list are \alert{not} the linear algebra operation.
\begin{defblock}An \alert{array} represents a vector or a matrix in linear algebra. It is often initialised from a list or another vector. Operations \pyth{+}, \pyth!*!, \pyth!/!, \pyth!-! are all \alert{elementwise}. The function \pyth!dot! or the operator \pyth!@! is used for the scalar product.
\end{defblock}

\end{frame}

\begin{frame}[fragile]{Vector usage}
\begin{example}
\begin{python}
vec = array([1., 3.]) # a vector in the plane
2*vec # array([2., 6.])
vec * vec # array([1., 9.])
vec/2 # array([0.5, 1.5])
norm(vec) # norm
vec @ vec # scalar product
\end{python}
\end{example}
\end{frame}

\begin{frame}[fragile]{Vectors are similar to lists}
\begin{itemize}
\item Access vectors via their indices
\begin{python}
v = array([1., 2., 3.])
v[0] # 1.
\end{python}
\item The \alert{length} of a vector is obtained by the function \pyth!len!.
\begin{python}
len(v) # 3
\end{python}
\item Parts of vectors using slices
\begin{python}
v[:2] # array([1., 2.])
\end{python}
\item Replace parts of vectors using slices
\begin{python}
v[:2] = [10, 20]
v # array([10., 20., 3.])
\end{python}
\end{itemize}
\end{frame}

\begin{frame}[fragile]{Vectors are not lists!}
Operations are not the same:
\begin{itemize}
\item Operations \pyth{+} and \pyth{*} are \alert{different}
\item More operations are defined: \pyth{-}, \pyth{/}
\item Many functions act elementwise on vectors: \pyth!exp!, \pyth!sin!, \pyth!sqrt!, etc.
\item Scalar product with \pyth!@!
\end{itemize}
\vfill
\begin{itemize}
\item Vectors have a \alert{fixed size}: no \pyth{append} method
\item Only \alert{one type} throughout the whole vector\\ \hfill (usually \code{float}, \code{complex}, \code{int} or \code{bool})
\item Vector slices are \alert{views}:
\begin{python}
v = array([1., 2., 3.])
v1 = v[:2] # v is array([1.,2.])
v1[0] = 0. # if v1 is changed ...
v # ... v is changed too: array([0., 2., 3.])
\end{python}
\end{itemize}
\end{frame}
\begin{frame}[fragile]{More examples}
\begin{python}
v1 = array([1., 2., 3.]) # don't forget the dots!
v2 = array([2, 0, 1.]) # one dot is enough
v1 + v2; v1 / v2; v1 - v2; v1 * v2

3*v1
3*v1 + 2*v2
\end{python}
\vfill
\begin{python}
v1 @ v2 # scalar product
cos(v1) # cosine, elementwise

# access
v1[0] # 1.
v1[0] = 10

# slices
v1[:2] # array([10., 2.])
v1[:2] = [0, 1] # now v1 == array([0.,1.,3.])
v1[:2] = [1,2,3] # error!
\end{python}
\end{frame}
\begin{frame}[fragile]{Creating Vectors \hfill (linspace)}
The \pyth!linspace! method is a convenient way to create equally spaced arrays.
\begin{python}
xs = linspace(0, 10, 200) # 200 points between 0 and 10
xs[0] # the first point is 0
xs[-1] # the last is 10
\end{python}
So for example the plot of the sine function between 0 and 10 will be obtained by:
\begin{python}
plot(xs, sin(xs))
\end{python}
\end{frame}

\begin{frame}[fragile]{Creating Vectors \hfill (zeros, ones, ...)}
Some handy methods to quickly create vectors:
\begin{description}
\item[zeros] \pyth!zeros(n)! creates a vector of size $n$ filled with zeros
\item[ones] \pyth!ones(n)! is the same filled with ones
\item[full] \pyth!full(n, value)! creates a vector of size $n$ filled with \pyth!value!
\item[rand] \pyth!rand(n)! creates a vector with {\tiny uniformly distributed} random numbers between 0 and 1
\item[empty] \pyth!empty(n)! creates an ``empty'' vector of size $n$ \\ \hfill (try it!)
\end{description}
\end{frame}

\begin{frame}[fragile]{Concatenating Vectors}
We need a means to \emph{concatenate} vectors. This is where the command \pyth!hstack! comes to help. \pyth!hstack([v1, v2,..., vn])! concatenates the vectors \pyth!v1!, \pyth!v2!, \ldots, \pyth!vn!.
\vfill
\begin{exampleblock}{Symplectic permutation}
We have a vector of size $2n$. We want to permute the first half with the second half of the vector with sign change:
\[(x_1, x_2,\ldots,x_n,x_{n+1},\ldots,x_{2n})\mapsto (-x_{n+1}, -x_{n+2},\ldots,-x_{2n},x_1,\ldots,x_n)\]
\begin{python}
def symp(v):
  n = len(v) // 2 # use the integer division //
  return hstack([-v[-n:], v[:n]])
\end{python}
\end{exampleblock}
\end{frame}

\begin{frame}[fragile]{Vectorized Functions}
Note that \alert{not all functions} may be applied on vectors. For instance this one:
\begin{python}
def const(x):
  return 1
\end{python}
We will see later how to automatically \emph{vectorize} a function so that it works componentwise on vectors.
\end{frame}


\begin{frame}[fragile]{Matrices as Lists of Lists}
\begin{defblock}
Matrices are represented by arrays of lists of \alert{rows}, which are lists as well.
\end{defblock}
\begin{python}
# the identity matrix in 2D
id = array([[1., 0.], [0., 1.]])
# Python allows this:
id = array([[1., 0.],
            [0., 1.]])
# which is more readable
\end{python}
\end{frame}

\begin{frame}[fragile]{Accessing Matrix Entries}
Matrix coefficients are accessed with \alert{two} indices:
\begin{python}
M = array([[1., 2.],[3.,4.]])
M[0,0] # first row, first column: 1.
M[-1,0] # last row, first column: 3.
\end{python}
\end{frame}

\begin{frame}[fragile]{Creating Matrices}
Some convenient methods to create matrices are:
\begin{description}
\item[eye] \pyth!eye(n)! is the identity matrix of size $n$
\item[zeros] \pyth!zeros([n,m])! fills an $n\times m$ matrix with zeros
\item[full] \pyth!full((n,m),a)! fills an $n\times m$ matrix with the value given by $a$
\item[rand] \pyth!rand(n,m)! is the same with random values
\item[empty] \pyth!empty([n,m])! same with ``empty'' values
\end{description}
\end{frame}

\begin{frame}[fragile]{Shape}
The \alert{shape} of a matrix is the tuple of its dimensions. The shape of an $n\times m$ matrix is \pyth!(n,m)!. It is given by the method \pyth{shape}:
\begin{python}
M = eye(3)
M.shape # (3, 3)

V = array([1., 2., 1., 4.])
V.shape # (4,) <- tuple with one element
\end{python}
\vfill
Tip:
\begin{itemize}
\item[] \pyth!zeros_like(A)! returns a matrix of the same shape as $A$ containing only zeros.
\item[] \pyth!ones_like, empty_like, full_like! does the same but corresponding to the commands \pyth!ones, empty, full!.
\end{itemize}
\vfill


\end{frame}

\begin{frame}[fragile]{Transpose}
The \alert{transpose} of a matrix $A_{ij}$ is a matrix $B$ such that
\[B_{ij} = A_{ji}\]

By transposing a matrix you \alert{switch} the two shape elements.

%You may \alert{switch} the two shape elements by \alert{transposing} the matrix.
%The \alert{transpose} of a matrix $A_{ij}$ is a matrix $B$ such that
%\[B_{ij} = A_{ji}\]
\begin{python}
A = ...
A.shape # 3,4

B = A.T # A transpose
B.shape # 4,3
\end{python}
\vfill
Note: $B$ is just a ``view'' of $A$. Altering $B$ changes $A$.
\end{frame}

\begin{frame}[fragile]{Matrix vector multiplication}
The mathematical concept of \emph{reduction}:
\[ \sum_{j} a_{ij} x_{j}\]
is translated in Python in the function \pyth!dot!:
\begin{python}
angle = pi/3
M = array([[cos(angle), -sin(angle)],
           [sin(angle), cos(angle)]])
v = array([1., 0.])
Y = M @ v # the product M v
\end{python}
% \pause
\begin{alertblock}{Elementwise vs. matrix multiplication}
The multiplication operator \pyth!*! is \alert{always} elementwise. It has nothing to do with the dot operation. \pyth!A*V! is a legal operation which will be explained later on.
\end{alertblock}
\end{frame}

\begin{frame}[fragile]{Dot product}
\begin{description}
\item[vector vector]  \hfill $s = \sum_ix_i y_i$ \hfill \pyth!s = x @ y!
\vfill
\item[matrix vector ]\hfill $ y_i =\sum_j A_{ij}x_j$ \hfill \pyth!y = A @ x!
 \vfill
\item[matrix matrix] \hfill $C_{ij} =\sum_{k} A_{ik}B_{kj} $ \hfill \pyth!C = A @ B!
 \vfill
\item[vector matrix] \hfill $y_j =\sum_i  x_i A_{ij}$ \hfill \pyth!y = x @ A!
\vfill
\end{description}
\end{frame}
\begin{frame}[fragile]{Solving a Linear System}
If $A$ is a matrix and $b$ is a vector you solve the linear equation
\[ A\cdot x = b\]
using \pyth!solve! which has the syntax \pyth!x = solve(A,b)!.

\begin{example}
We want to solve
\[\left\{\begin{aligned}x_1 &+ 2x_2 &= 1 \\ 3x_1 &+ 4x_2 &= 4\end{aligned}\right.\]
\begin{python}
from scipy.linalg import solve

A = array([[1., 2.],
           [3., 4.]])
b = array([1., 4.])
x = solve(A,b)
A @ x  # should be almost b
\end{python}
\end{example}
\end{frame}

%\subsection{Matrix Slices}

\begin{frame}[fragile]{Slices}
Slices are similar to that of lists and vectors \alert{except} that there are now \alert{two dimensions}.
\begin{description}
\item[\texttt{M[i,:]}] a \emph{vector} filled by the row $i$ of $M$
\item[\texttt{M[:,j]}] a \emph{vector} filled by the column $j$ of $M$
\item[\texttt{M[2:4,:]}] slice \pyth!2:4! on the rows only
\item[\texttt{M[2:4,1:4]}] slice on rows and columns
\end{description}

\begin{alertblock}{Omitting a dimension}
If you omit an index or a slice, \scipy \hspace{.2em} assumes you are taking \alert{rows only}.
\begin{description}
\item[\texttt{M[3]}] is the third row of $M$
\item[\texttt{M[1:3]}] is a matrix with the second and third rows of $M$.
\end{description}
\end{alertblock}
\end{frame}

\begin{frame}[fragile]{Altering a Matrix}
You may alter a matrix using slices or direct access.
\begin{itemize}
\item \pyth!M[2,3] = 2.!
\item \pyth!M[2,:] = <a vector>!
\item \pyth!M[1:3,:] = <a matrix>!
\item \pyth!M[1:4,2:5] = <a matrix>!
\end{itemize}
The matrices and vectors above \alert{must have the right size} to ``fit'' in the matrix $M$.
\end{frame}




\begin{frame}[fragile]{Some technical terms}
\vfill
\begin{itemize}
 \item[Rank] The rank of an array is the number of indices used to identify an element.
A matrix has rank 2, a vector has rank 1, a scalar has rank 0.\\ \emph{Do not confound this term with the
rank used in linear algebra.}
\item[Shape] The shape is a tuple with the dimensions of the array: A $2 \times 3$ matrix is
represented in Python by an array with shape $(2,3)$.
\item[Length] The number of elements in an array. An array with shape $(2,3)$ has length $6$.
An array with shape $(2,)$ (representing a vector) has length $2$.\\
\emph{Do not confound this term with the term length used in linear algebra.}
\end{itemize}
\vfill
\end{frame}

\begin{frame}[fragile]{Rank of matrix slices}
When slicing the rank of the result object is as follows:
\begin{center}
\begin{tabular}{l|l|l}
access & rank & kind\\
\hline\hline
index,index & 0 & scalar \\\hline
slice,index & 1 & vector \\\hline
slice,slice & 2 & matrix\\
\end{tabular}
\end{center}
% \pause
\begin{example}
\begin{columns}
\begin{column}{4cm}
\begin{overlayarea}{4cm}{3.5cm}
\begin{onlyenv}<2>
\begin{tikzpicture}
\matrix[nodes={draw,style=empty},minimum size= .8cm,ampersand replacement=\&]
{
\node{0} ; \& \node[style=m]{1}; \& \node[style=m]{2};  \& \node{3}; \\
\node{4}; \& \node[style=m]{5}; \& \node[style=m]{6}; \& \node{7}; \\
\node{8}; \& \node{9}; \& \node{10}; \& \node{11}; \\
};
\end{tikzpicture}
\end{onlyenv}
\begin{onlyenv}<3>
\begin{tikzpicture}
\matrix[nodes={draw,style=empty},minimum size= .8cm,ampersand replacement=\&]
{
\node{0} ; \& \node{1}; \& \node{2};  \& \node{3}; \\
\node[style=v]{4}; \& \node[style=v]{5}; \& \node[style=v]{6}; \& \node[style=v]{7}; \\
\node{8}; \& \node{9}; \& \node{10}; \& \node{11}; \\
};
\end{tikzpicture}
\end{onlyenv}
\begin{onlyenv}<4>
\begin{tikzpicture}
\matrix[nodes={draw,style=empty},minimum size= .8cm,ampersand replacement=\&]
{
\node{0} ; \& \node{1}; \& \node{2};  \& \node{3}; \\
\node{4}; \& \node[style=s]{5}; \& \node{6}; \& \node{7}; \\
\node{8}; \& \node{9}; \& \node{10}; \& \node{11}; \\
};
\end{tikzpicture}
\end{onlyenv}
\begin{onlyenv}<5>
\begin{tikzpicture}
\matrix[nodes={draw,style=empty},minimum size= .8cm,ampersand replacement=\&]
{
\node{0} ; \& \node{1}; \& \node{2};  \& \node{3}; \\
\node [style=m]{4}; \& \node [style=m]{5}; \& \node [style=m]{6}; \& \node [style=m]{7}; \\
\node{8}; \& \node{9}; \& \node{10}; \& \node{11}; \\
};
\end{tikzpicture}
\end{onlyenv}
\begin{onlyenv}<6>
\begin{tikzpicture}
\matrix[nodes={draw,style=empty},minimum size= .8cm,ampersand replacement=\&]
{
\node{0} ; \& \node{1}; \& \node{2};  \& \node{3}; \\
\node{4}; \& \node[style=m]{5}; \& \node{6}; \& \node{7}; \\
\node{8}; \& \node{9}; \& \node{10}; \& \node{11}; \\
};
\end{tikzpicture}
\end{onlyenv}
\end{overlayarea}
\end{column}
\begin{column}{7cm}
\begin{tabular}{l|l|l|l}
\hline
\rowcolor{white}access & shape & rank & kind \\
\hline\hline
\rowcolor{matcolour}{\pyth!M[:2,1:-1]!} & $(2,2)$     & 2 & matrix\\ \hline
\rowcolor{veccolour}{\pyth!M[1,:]!}     &  $4$        & 1 & vector\\ \hline
\rowcolor{scalcolour}{\pyth!M[1,1]!}    & $\emptyset$ & 0 & scalar\\ \hline
\rowcolor{matcolour}{\pyth!M[1:2,:]!}   & $(1,4)$     & 2 & matrix\\ \hline
\rowcolor{matcolour}{\pyth!M[1:2,1:2]!} & $(1,1)$     & 2 & matrix\\
\end{tabular}
\end{column}
\end{columns}
\end{example}
\end{frame}
\begin{frame}[fragile]{Reshaping}
From a given tensor (vector or matrix) one may obtain another tensor by \alert{reshaping}.
\end{frame}


\begin{frame}[fragile]{Reshaping Example}
\begin{columns}
\begin{column}{5cm}
\begin{onlyenv}<1>
\begin{python}
A = arange(6)
\end{python}
\end{onlyenv}
\begin{onlyenv}<2>
\begin{python}
A.reshape(1,6)
\end{python}
\end{onlyenv}
\begin{onlyenv}<3>
\begin{python}
A.reshape(6,1)
\end{python}
\end{onlyenv}
\begin{onlyenv}<4>
\begin{python}
A.reshape(2,3)
\end{python}
\end{onlyenv}
\begin{onlyenv}<5>
\begin{python}
A.reshape(3,2)
\end{python}
\end{onlyenv}
\end{column}
\begin{column}{5cm}
\begin{onlyenv}<1>
\begin{tikzpicture}
\matrix[nodes={style=empty},minimum size= .8cm,ampersand replacement=\&]
{
\node{0} ; \& \node{1}; \& \node{2};  \& \node{3}; \& \node{4}; \& \node{5}; \\
};
\end{tikzpicture}
\end{onlyenv}
\begin{onlyenv}<2>
\begin{tikzpicture}
\matrix[nodes={draw,style=empty},minimum size= .8cm,ampersand replacement=\&]
{
\node{0} ; \& \node{1}; \& \node{2};  \& \node{3}; \& \node{4}; \& \node{5}; \\
};
\end{tikzpicture}
\end{onlyenv}
\begin{onlyenv}<3>
\begin{tikzpicture}
\matrix[nodes={draw,style=empty},minimum size= .8cm,ampersand replacement=\&]
{
\node{0}; \\  \node{1}; \\  \node{2}; \\ \node{3}; \\  \node{4}; \\ \node{5}; \\ \\
};
\end{tikzpicture}
\end{onlyenv}
\begin{onlyenv}<4>
\begin{tikzpicture}
\matrix[nodes={draw,style=empty},minimum size= .8cm,ampersand replacement=\&]
{
\node{0} ; \& \node{1}; \& \node{2}; \\   \node{3}; \& \node{4}; \& \node{5}; \\
};
\end{tikzpicture}
\end{onlyenv}
\begin{onlyenv}<5>
\begin{tikzpicture}
\matrix[nodes={draw,style=empty},minimum size= .8cm,ampersand replacement=\&]
{
\node{0} ; \& \node{1}; \\ \node{2};  \& \node{3}; \\ \node{4}; \& \node{5}; \\
};
\end{tikzpicture}
\end{onlyenv}
\end{column}
\end{columns}
\end{frame}
\begin{frame}[fragile]{Reshaping Trick}
Note that Python can \alert{guess one of the new dimensions}. Just give a negative integer for the dimension to be guessed:
\begin{python}
A = arange(12) # a vector of length 12

A.reshape(1,-1) # row matrix
A.reshape(-1,1) # column matrix

A.reshape(3,-1) # 3,4 matrix
A.reshape(-1,4) # same
\end{python}
\end{frame}

%\subsection{Building matrices}

\begin{frame}[fragile]{Building Matrices}
\begin{itemize}
\item Piling vectors
\item Stacking vectors
\item Stacking column matrices
\end{itemize}
% \pause
The universal method to build matrices is \pyth!concatenate!. This function is called by several convenient functions
\begin{itemize}
\item \pyth!hstack! to stack matrices \alert{horizontally}
\item \pyth!vstack! to stack matrices \alert{vertically}
\item \pyth!column_stack! to stack \alert{vectors} in columns
\end{itemize}
\end{frame}

\begin{frame}[fragile]{Stacking Vectors}
\begin{columns}
\begin{column}{5cm}
\begin{python}
v1 = array([1,2])
v2 = array([3,4])
\end{python}
\begin{python}
vstack([v1,v2])
\end{python}
\begin{python}
column_stack([v1,v2])
\end{python}
\end{column}
\begin{column}{5cm}
\begin{tikzpicture}
\matrix[nodes={draw,style=empty},minimum size= .8cm,ampersand replacement=\&]
{
\node[style=one]{1} ; \& \node[style=one]{2}; \\ \node[style=two]{3};  \& \node[style=two]{4}; \\
};
\end{tikzpicture}
\begin{tikzpicture}
\matrix[nodes={draw,style=empty},minimum size= .8cm,ampersand replacement=\&]
{
\node[style=one]{1} ; \& \node[style=two]{3}; \\ \node[style=one]{2};  \& \node[style=two]{4}; \\
};
\end{tikzpicture}
\end{column}
\end{columns}
\end{frame}

%\subsection{Array Methods}

\begin{frame}[fragile]{sum, max, min}
\begin{columns}
\begin{column}{5cm}
You may perform a number of operations on arrays, either on the whole array, or column-wise or row-wise. The most common are
\begin{itemize}
\item\alert{\textit{max}}
\item \alert{\textit{min}}
\item \alert{\textit{sum}}
\end{itemize}
\end{column}
\begin{column}{5cm}
\begin{example}
\begin{tikzpicture}
\matrix[nodes={draw,style=empty},minimum size= .8cm,ampersand replacement=\&]
{
\node{1} ; \& \node{2}; \& \node{3};  \& \node{4}; \\
\node{5} ; \& \node{6}; \& \node{7};  \& \node{8}; \\
};
\end{tikzpicture}
\begin{python}
A.sum()
\end{python}
36
\end{example}
\end{column}
\end{columns}
\end{frame}
\begin{frame}[fragile]{sum, max, min \hfill axis parameter}
\begin{columns}
\begin{column}{5cm}
\begin{tikzpicture}
\matrix[nodes={draw,style=empty},minimum size= .8cm,ampersand replacement=\&]
{
\node[style=one]{1} ; \& \node[style=two]{2}; \& \node[style=one]{3};  \& \node[style=two]{4}; \\
\node[style=one]{5} ; \& \node[style=two]{6}; \& \node[style=one]{7};  \& \node[style=two]{8}; \\
};
\end{tikzpicture}
\begin{python}
A.sum(axis=0)
\end{python}
The result is a \alert{vector}
\begin{tikzpicture}
\matrix[nodes={style=empty},minimum size= .8cm,ampersand replacement=\&]
{
\node[style=one]{6} ; \& \node[style=two]{8}; \& \node[style=one]{10};  \& \node[style=two]{12}; \\
};
\end{tikzpicture}

\end{column}
\begin{column}{5cm}
\begin{tikzpicture}
\matrix[nodes={draw,style=empty},minimum size= .8cm,ampersand replacement=\&]
{
\node[style=one]{1} ; \& \node [style=one]{2}; \& \node [style=one]{3};  \& \node [style=one]{4}; \\
\node [style=two]{5} ; \& \node[style=two]{6}; \& \node [style=two]{7};  \& \node[style=two]{8}; \\
};
\end{tikzpicture}
\begin{python}
A.sum(axis=1)
\end{python}
The result is a \alert{vector}
\begin{tikzpicture}
\matrix[nodes={style=empty},minimum size= .8cm,ampersand replacement=\&]
{
\node[style=one]{10} ; \&  \node[style=two]{26}; \\
};
\end{tikzpicture}
\end{column}
\end{columns}
\end{frame}
\begin{frame}[fragile]{Boolean Arrays}
One may use \alert{Boolean arrays} to create a ``template'' for modifying another array:\\
\vfill
\begin{python}
B = array([[True, False],
           [False, True]])  # the template array
M = array([[2, 3],
           [1, 4]])         # the other array
M[B] = 0    #  using the template
M # [[0, 3], [1, 0]]
M[B] = 10, 20
M # [[10, 3], [1, 20]]
\end{python}
\end{frame}

\begin{frame}[fragile]{Creating Boolean Arrays}
It might be just as tedious to create the boolean array by hand than to change the array directly. There are however many methods to create Boolean arrays.
% \pause

Any \alert{logical operator} will create a Boolean array instead of a Boolean.

\begin{python}
M = array([[2, 3], [1, 4]])
M > 2 # array([[False, True], [False, True]])
M == 0 # array([[False, False], [False, False]])
N = array([[2, 3], [0, 0]])
M == N # array([[True, True], [False, False]])
...
\end{python}

This allows the elegant syntax:
\begin{python}
M[M>2] = 0
# all the elements > 2 are replaced by 0
\end{python}
\end{frame}

\subsection{Comparing Arrays}

\begin{frame}[fragile]{Comparing Arrays}
Note that because array comparison create Boolean arrays, \alert{one cannot compare arrays directly}.


The solution is to use the methods \pyth!all! and \pyth!any!:
\begin{python}
A = array([[1,2],[3,4]])
B = array([[1,2],[3,3]])
A == B # creates array([[True, True], [True, False]])
(A == B).all() # False
(A != B).any() # True
\end{python}
\end{frame}

\begin{frame}[fragile]{Boolean Operations}
For the same reason as before you
\alert{cannot} use \pyth|and|, \pyth|or| nor \pyth!not! on Boolean arrays! Use the following replacement operators instead:

\newcommand{\mytilde}{\raise.17ex\hbox{$\scriptstyle\mathtt{\sim}$}}
\begin{center}
\rowcolors[]{1}{blue!20}{blue!10}
\begin{tabular}{c|c}
\hline
logic operator & replacement for Boolean arrays\\ \hline
A and B & A \& B \\ \hline
A or B & A | B\\ \hline
not A & \mytilde A \\ \hline
\end{tabular}
\end{center}

\begin{python}
a = array([True, True, False, False])
b = array([True, False, True, False])

a and b # error!
a & b # array([True, False, False, False])
a | b # array([True, Trues, True, False])
~a # array([False, False, True, True])
\end{python}
\end{frame}
\begin{frame}[fragile]{Universal Functions}
\begin{defblock}
A universal function (or ufunc) is a function that operates on arrays in an element-by-element fashion. That is, a ufunc is a “vectorized” wrapper for a function that takes a fixed number of scalar inputs and produces a fixed number of scalar outputs.
\end{defblock}
\vfill
Examples:
\begin{python}
from scipy import *
sin
cos
exp
\end{python}
\end{frame}
\begin{frame}[fragile]{Vectorized Functions}
Non-universal functions can be wrapped to behave like universal funtions.\\
This is done by the command \pyth!vectorize!.\\
\begin{python}
# A non-universal function
def heaviside(x):
   if x >= 0:
      return 1.
   else:
      return 0.

heaviside(linspace(-1,1,100))  # returns an error

vheaviside=vectorize(heaviside) # a new function
vheaviside(linspace(-1,1,100)) # does the job
\end{python}
\end{frame}
\begin{frame}[fragile]{Vectorized Functions \hfill (Cont.)}
\begin{python}
xvals=linspace(-1,1,100)
plot(xvals,vectorize(heaviside)(xvals))
axis([-1.5,1.5,-0.5,1.5])
\end{python}
\end{frame}

\end{document}
